\documentclass[12pt]{article}
\makeatletter
\def\input@path{{../../}}
\makeatother
\usepackage{sbc-template}
\makeatletter
\renewcommand{\inst}[1]{}
\makeatother
\usepackage{graphicx,url}
\usepackage[utf8]{inputenc}
\usepackage[T1]{fontenc}
\usepackage[brazil]{babel}
\usepackage{longtable,booktabs,array}
\providecommand{\tightlist}{%%
  \setlength{\itemsep}{0pt}\setlength{\parskip}{0pt}}
\sloppy
\title{De Reservoir Computing classico a Quantum Reservoir Computing: a ponte conceitual}
\author{}
\address{}
\begin{document}

\maketitle

\begin{resumo}
Este artigo prepara a transicao intelectual da serie. Ate aqui, trabalhamos com Reservoir Computing como um paradigma implementado em software classico. Agora precisamos responder a pergunta que torna o artigo final possivel: o que muda, o que permanece e por que Quantum Reservoir Computing pode ser lido como uma extensao natural de RC? A resposta proposta neste texto e que RC deve ser entendido antes de tudo como um paradigma de computacao por dinamica rica mais readout simples. Quando essa chave conceitual esta clara, reservatorios fisicos e reservatorios quanticos deixam de parecer exoticos e passam a ser variacoes de uma mesma ideia estrutural. O artigo apresenta esse paralelo, discute limites praticos de QRC e define como o caso Favorita e adaptado para um pipeline quantico controlado.
\end{resumo}

\section{1. Introducao}\label{1-introducao}

Se o leitor chegou ate aqui, ja conhece o problema de negocio, ja viu o pipeline classico e ja aprendeu a interpretar RC sem mistica. O desafio agora e evitar um salto conceitual brusco. Em muitos textos sobre QRC, o quantico entra cedo demais e o leitor perde o fio do argumento. Nesta serie, o caminho e o inverso: primeiro consolidamos o paradigma, depois trocamos o suporte dinamico.

Essa troca so faz sentido quando a pergunta central esta bem posta:

\begin{quote}
O que, em RC, e essencial? E o que, em QRC, e apenas uma nova realizacao desse essencial?
\end{quote}

\section{2. RC como paradigma geral}\label{2-rc-como-paradigma-geral}

Reservoir Computing pode ser resumido em uma estrutura de tres blocos:

\begin{enumerate}
\def\labelenumi{\arabic{enumi}.}
\tightlist
\item
  um mecanismo de codificacao da entrada;
\item
  um sistema dinamico rico que transforma a entrada em estado;
\item
  um readout simples treinado sobre observaveis desse estado.
\end{enumerate}

Essa descricao e poderosa porque e mais geral do que uma Echo State Network. Ela nao exige que o reservatorio seja digital, neural ou mesmo simulavel da forma tradicional. O que importa e a combinacao entre:

\begin{itemize}
\tightlist
\item
  riqueza dinamica;
\item
  resposta dependente do historico;
\item
  leitura treinavel e relativamente simples.
\end{itemize}

Quando entendemos RC assim, reservatorios fisicos deixam de ser excecao. Eles passam a ser instancias concretas do paradigma.

\section{3. Reservatorios fisicos}\label{3-reservatorios-fisicos}

A literatura de physical reservoir computing mostra que o reservatorio pode ser realizado em sistemas opticos, eletricos, mecanicos, analogicos e outros meios fisicos. O motivo e direto: nao precisamos treinar toda a dinamica, apenas explorar uma dinamica suficientemente rica e observavel.

Essa ideia tem forte valor pedagogico para a ponte com QRC:

\begin{itemize}
\tightlist
\item
  o reservatorio nao precisa ser um grafo neural tradicional;
\item
  o foco deixa de ser arquitetura de treino e passa a ser dinamica observavel;
\item
  a fronteira entre modelo e meio fisico fica mais interessante.
\end{itemize}

O passo para o quantico, entao, nao e abandonar RC. E perguntar se um sistema quantico pode cumprir o mesmo papel de reservatorio.

\section{4. O salto para o quantico}\label{4-o-salto-para-o-quantico}

No QRC, os blocos conceituais permanecem:

\begin{itemize}
\tightlist
\item
  ha uma entrada;
\item
  ha uma dinamica interna;
\item
  ha observaveis;
\item
  ha um readout treinado.
\end{itemize}

O que muda e o suporte dinamico.

Em vez de um estado vetorial classico atualizado por uma regra recorrente, temos:

\begin{itemize}
\tightlist
\item
  um estado quantico;
\item
  uma evolucao controlada por portas ou Hamiltonianos;
\item
  observaveis obtidos por medidas ou valores esperados;
\item
  um readout classico treinado sobre essas observacoes.
\end{itemize}

Em outras palavras, QRC nao elimina a logica de RC. Ele a reimplementa em um espaco de estados diferente.

\section{5. Paralelo direto entre RC e QRC}\label{5-paralelo-direto-entre-rc-e-qrc}

O quadro abaixo resume a continuidade entre os dois paradigmas.

\begin{longtable}[]{@{}lll@{}}
\toprule\noalign{}
Elemento & RC classico & QRC \\
\midrule\noalign{}
\endhead
\bottomrule\noalign{}
\endlastfoot
Entrada & vetor classico \texttt{u\_t} & vetor classico codificado em portas ou parametros \\
Estado & vetor dinamico \texttt{x\_t} & estado quantico do circuito ou sistema \\
Dinamica & recorrencia classica fixa & evolucao quantica fixa ou controlada \\
Observacao & estado interno ou funcoes dele & valores esperados de observaveis \\
Readout & regressao ou classificador simples & regressao ou classificador classico simples \\
\end{longtable}

Esse paralelo e o coracao da ponte conceitual da serie. Sem ele, QRC parece um tema separado. Com ele, QRC aparece como uma extensao natural da mesma pergunta: qual dinamica produz uma representacao temporal util?

\section{6. O que se ganha e o que se perde}\label{6-o-que-se-ganha-e-o-que-se-perde}

\subsection{6.1 O que se ganha}\label{61-o-que-se-ganha}

\begin{itemize}
\tightlist
\item
  um espaco de estados potencialmente muito rico;
\item
  dinamicas que nao tem analogia imediata em reservatorios classicos simples;
\item
  uma agenda de pesquisa que conecta machine learning, simulacao e sistemas quanticos.
\end{itemize}

\subsection{6.2 O que se perde}\label{62-o-que-se-perde}

\begin{itemize}
\tightlist
\item
  simplicidade computacional;
\item
  escalabilidade imediata;
\item
  facilidade de interpretacao;
\item
  comparabilidade direta com pipelines classicos em datasets grandes.
\end{itemize}

Essa assimetria precisa ser explicitada. O artigo final nao existe para "provar que o quantico vence". Ele existe para mostrar como montar um experimento inicial de QRC que seja metodologicamente honesto e pedagogicamente instrutivo.

\section{7. Como o projeto adapta o caso Favorita para QRC}\label{7-como-o-projeto-adapta-o-caso-favorita-para-qrc}

No repositorio desta serie, o pipeline quantico foi desenhado para preservar o mesmo recorte do benchmark:

\begin{itemize}
\tightlist
\item
  dataset: Favorita \texttt{train.csv}
\item
  serie: \texttt{store\_nbr\ =\ 1}, \texttt{family\ =\ BEVERAGES}
\item
  teste: ultimos 90 dias
\end{itemize}

Ao mesmo tempo, o projeto reconhece o custo da dinamica quantica e faz escolhas controladas:

\begin{itemize}
\tightlist
\item
  usa um vetor quantico de entrada reduzido;
\item
  usa um readout linear classico;
\item
  executa tuning focal em \texttt{n\_qubits} e \texttt{window};
\item
  compara o QRC com o mesmo benchmark classico ja estabelecido.
\end{itemize}

Mais especificamente, o pipeline quantico do projeto codifica quatro sinais:

\begin{itemize}
\tightlist
\item
  demanda previa normalizada;
\item
  promocao normalizada;
\item
  \texttt{dow\_sin};
\item
  \texttt{dow\_cos}.
\end{itemize}

Essa reducao nao e um defeito escondido. E uma decisao metodologica explicita para manter o experimento viavel sem quebrar o fio comparativo da serie.

\section{8. Por que o benchmark anterior precisa ser herdado}\label{8-por-que-o-benchmark-anterior-precisa-ser-herdado}

Sem o benchmark do artigo 5, QRC poderia ser apresentado em um ambiente artificialmente favoravel. Ao herdar o mesmo recorte e as mesmas metricas, o artigo final se protege de uma armadilha comum: comparar o modelo quantico a um problema diferente daquele usado pelos classicos.

Aqui, a comparacao continua sendo:

\begin{itemize}
\tightlist
\item
  mesmo problema de negocio;
\item
  mesmo alvo;
\item
  mesmo horizonte;
\item
  mesmas metricas;
\item
  nova forma de reservatorio.
\end{itemize}

Essa disciplina nao apenas melhora a ciencia. Ela melhora o ensino.

\section{9. Expectativa correta para o artigo final}\label{9-expectativa-correta-para-o-artigo-final}

O leitor deve chegar ao artigo 7 com uma expectativa equilibrada:

\begin{itemize}
\tightlist
\item
  o QRC deve funcionar de ponta a ponta;
\item
  o QRC deve ser comparado com honestidade;
\item
  o QRC nao precisa vencer os classicos para ser cientificamente valioso;
\item
  o importante e entender o pipeline, os trade-offs e o tipo de pergunta que o metodo permite fazer.
\end{itemize}

Essa expectativa e central para uma obra 101. O objetivo nao e hype. O objetivo e formar criterio.

\section{10. Conclusao}\label{10-conclusao}

RC e, antes de tudo, um paradigma de computacao por dinamica rica e leitura simples. Quando essa ideia esta clara, QRC pode ser introduzido sem salto arbitrario: trocamos o tipo de dinamica, preservamos a logica metodologica.

Com a ponte conceitual estabelecida, estamos prontos para o artigo final da serie: um estudo didatico de QRC aplicado ao problema de previsao de demanda no Favorita.

\section{Entregaveis associados no repositorio}\label{entregaveis-associados-no-repositorio}

\begin{itemize}
\tightlist
\item
  implementacao do QRC: \texttt{code/qrc/model.py}
\item
  resultados do QRC: \texttt{code/qrc/RESULTS.md}
\item
  tuning de qubits e janela: \texttt{code/qrc/experiments/TUNING.md}
\end{itemize}

\section{Referencias}\label{referencias}

\begin{enumerate}
\def\labelenumi{\arabic{enumi}.}
\tightlist
\item
  Tanaka, G. et al. Recent advances in physical reservoir computing: A review.
\item
  Fujii, K., Nakajima, K. Harnessing disordered-ensemble quantum dynamics for machine learning.
\item
  Jaeger, H. The "echo state" approach to analysing and training recurrent neural networks.
\end{enumerate}

\section*{Resultados Computacionais Associados}

Os resultados computacionais associados a este artigo estao organizados na subpasta \texttt{computational\_results\_20260402\_222902/}, localizada dentro desta pasta \LaTeX. Esse diretorio contem o arquivo \texttt{REPORT.md}, tabelas em CSV e imagens comparativas apropriadas ao escopo do artigo.

\end{document}
